% Generated by analysis/scripts/export_ten_cases.py
\subsection*{统一结果卡}
\begingroup\small
\noindent\begin{tabularx}{\textwidth}{@{}p{24mm}X@{}}\toprule
技术/模式 & HfO$_2$ FeRAM；HZO 1T1C二状态数字路径；已有读码跨8输入位复用\\
逻辑粒度/聚合 & B\_S=128 Byte；B\_R=16 Byte。128bit完整局部行＝16 Byte；1024行事务覆盖矩阵；每次真实破坏读均整行恢复\\
资源/相对R0 & 32片×128 SA/恢复驱动；已有4096bit读码寄存器，新增权重锁存为0；128目标写驱动\\
服务口径 & 真实破坏读恢复已含raw streaming；无另列周期刷新\\
主导因素/选择 & 32次真实读恢复＋256数字轮；每极性14/50/100 ns为保守工程预算，14 ns原值是write latency\\
关键对照 & 旧256次重读仅为调度对照；不将主动弃码归为介质必需恢复\\
\bottomrule\end{tabularx}
\vspace{3mm}
\begin{center}\begin{tabular}{@{}lrrr@{}}\toprule
成对条件情景 & $\rho$ (MB/s) & $\tau$ (MB/s) & $\mathrm{RI}^{*}$\\\midrule
短预算 & 68.82 & 307.7 & 0.2237\\
参考（推荐） & 26.61 & 106.7 & 0.2495\\
长预算 & 12.88 & 53.33 & 0.2414\\
\bottomrule\end{tabular}\end{center}
\noindent 参考原始占用：streaming 4810 ns；resident 150 ns。
\par\noindent 范围为有限成对条件情景极值，不是物理不确定性包络。1 MB=$10^6$ Byte；整矩阵16384 Byte=16 KiB。源定位见本例证据笔记；公共基线shared\_baseline。
\endgroup
\clearpage
